\documentclass[10pt,a4paper]{amsart}
\usepackage[margin=19mm]{geometry}
\usepackage{amsmath,amssymb,mathtools}
\usepackage[scr=rsfs,cal=euler]{mathalfa}
\usepackage{xfrac}
\usepackage[unicode,hidelinks,bookmarks=false]{hyperref}
\newcommand{\ie}{i.e.\ }
\newcommand{\cf}{cf.\ }
\newcommand{\resp}{respectively\ }
\newcommand{\Subsection}{\subsection}
\providecommand{\nobreakdash}{\nobreak\hbox{-}}
\setlength{\emergencystretch}{2em}
\multlinegap0pt
\usepackage{amssymb}
\usepackage{enumerate}
\usepackage{xcolor}
\newtheorem{thm}{Theorem}
\newcommand{\R}{\mathbb R}
\newcommand{\s}{\mathbb S}
\title{A constant rank theorem\\ for linear elliptic equations on the sphere\\ with applications to the mixed Christoffel problem}
\author{A. Colesanti, M. Focardi, P. Guan, P. Salani}
\date{Source: arXiv:2512.08670v2 (CC BY 4.0)}
\begin{document}
\maketitle

\noindent\emph{Source and licence.} A constant rank theorem\\ for linear elliptic equations on the sphere\\ with applications to the mixed Christoffel problem. Version arXiv:2512.08670v2 (CC BY 4.0), distributed by arXiv under the Creative Commons Attribution 4.0 International licence, http://creativecommons.org/licenses/by/4.0/, as recorded in the arXiv licence metadata for this exact version and shown on the abstract page https://arxiv.org/abs/2512.08670v2. Full licence notice: \texttt{licenses/arxiv-2512.08670-CC-BY-4.0.txt}.\par\smallskip

\noindent\textit{Editorial scope.} Counted unit: the article's thm 3dim (source0.tex lines 473--482). The article's complete original proof of it is delivered with it (source0.tex lines 670--705, one \texttt{proof} environment of 36 lines, which the article places away from the statement). No mathematics is written, repaired or re-derived: the statement and the proof are the article's own lines, in the article's own notation, and no retained line is substituted or re-flowed.  Retained uncounted as the source-local context the proof needs: lines 236--238; lines 282--284; lines 287--289; lines 291--293; lines 331--333; lines 349--366; lines 446--456.  Nothing was omitted from any retained span, so the package carries no omission record.  No retained line contains a citation, so the wrapper carries no bibliography and no bibliography entry.  No retained line contains a figure environment, an image-inclusion command or drawing code, so no image is delivered and the package declares no asset dependency.  Editorial formatting only, in addition: an AMS \texttt{amsart} preamble that restates the article's own declarations and macros used by the retained lines, namely the environments \texttt{thm} on separate counters, together with the standard packages those lines need.  The retained environments print in this document as Theorem 1 in order of appearance, whereas the article prints the same items with the numbering of its own full document.  The complete native source of the article as distributed in the arXiv e-print for this exact version is bundled unchanged as \texttt{sources/190678/source0.tex}.
\begin{equation}\label{GM condition}
\mbox{the $1$-homogeneous extension of $\dfrac 1f$ to $\R^{n+1}$, is convex.}    
\end{equation} 

\begin{equation}\label{cmc}
c_{ij}=\frac{\partial \mathscr{D}(M_1,\dots,M_{n-1},M)}{\partial m_{ij}},\quad\forall\, i,j=1,\dots,n.
\end{equation}

\begin{equation}\label{Prob-C}
\mathscr{D}(W_1,\cdots, W_{n-1}, W)=f, \quad \mbox{on $\s^n$},
\end{equation}

\begin{equation}\label{the equation}
\sum_{\alpha,\beta=1}^n a_{\alpha\beta}(x)(u_{\alpha\beta}(x)+\delta_{\alpha\beta}u(x))=f(x)
\end{equation}

\begin{equation}\label{cond-1 2} 
\left(\frac{a_{qq}}f\delta_{ij}+ \nabla_i\nabla_j\left(\frac{a_{qq}}{f}\right)\right)_{i,j=1,2}\ge 0.
\end{equation}

\begin{thm}\label{main theorem n=2 rephrased} Let $A(x)=(a_{ij}(x))_{i,j=1,2}\in C^{2,\gamma}(D)$ be a positive $(2,0)$-tensor defined in an open connected set $D\subset\s^2$, and let $f\in C^{2, \gamma}(D)$ for some $\gamma>0$, with $f>0$ on $D$. Let $M$ be the $1$-homogeneous extension of the tensor
$$
\frac{A}{f}
$$
to the cone 
$$
C=\left\{x\in\R^3\colon\frac{x}{|x|}\in D\right\}.
$$
Assume that  
\begin{equation}\label{the condition}
M\mbox{ is locally convex in $C$}.
\end{equation}
Let $u\in C^2(\s^2)$ be a solution of equation \eqref{the equation}. If  
\[
W=(u_{ij}+u\delta_{ij})_{i,j=1,2}\ge 0
\] 
in $D$, then $W$ has constant rank in $D$. In particular, if $D=\s^2$ then $W$ has full rank in $\s^2$. 
\end{thm}

\begin{thm}\label{MC 3dim 1} Let $\Omega_0\subset\R^3$ be a convex body with boundary of class $C^{4,\gamma}$, for some $\gamma>0$, and with positive Gauss curvature. Let $h$ be the support function of $\Omega_0$. Let $f\in C^{2,\gamma}(\s^2)$, with $f>0$ on $\s^2$, be such that condition \eqref{GM condition} holds. Assume that, for every $x\in\s^2$, and for every tangential direction $\alpha$ to $\s^2$ at $x$, the following inequality holds: 
\begin{equation}\label{new form dim two statement}
\left(\frac{h_{\alpha\alpha}(x)+h(x)}{f(x)}\right)_{\alpha\alpha}+\frac{h_{\alpha\alpha}(x)+h(x)}{f(x)}\ge0,
%\quad\forall\, x\in\s^2,\, \forall\, \alpha,
\end{equation}
where sub-indices denote covariant partial derivatives with respect to the direction $\alpha$. Then, there exists a convex body $\Omega\in C^{2,+}$ such that
$$
S(\Omega_0,\Omega,\cdot)=\mu(\cdot)\,,
$$
where $\mu$ is the measure such that $d\mu(x)=f(x)dx$.
\end{thm}

\begin{thm}\label{MC 3dim} Let $\Omega_0\subset\R^3$ and $f$ be as in Theorem \ref{MC 3dim 1}, and let $W_0$ be the inverse Weingarten map of $\Omega_0$. If $f$ verifies condition \eqref{GM condition}, and the $1$-homogeneous extension of the tensor 
$$
\frac{W_0}{f}
$$
is convex in $\R^3$, then there exists a convex body $\Omega\in C^{2,+}$ such that
$$
S(\Omega_0,\Omega,\cdot)=\mu(\cdot),
$$
where $\mu$ is the Borel measure on $\s^2$ such that $d\mu=f dx$.
\end{thm}

\begin{proof}[Proof of Theorem \ref{MC 3dim}] 
By \eqref{Prob-C}, the tensor $A=(a_{ij})_{i,j=1,2}$ appearing in \eqref{the equation} is given by
$$
A=C(W_0),
$$
where $C(W_0)$ is the co-factor matrix of $W_0$. Let
$$
W_0=
\left(
\begin{array}{cc}
b_{11}     &b_{12}  \\
b_{21}     &b_{22} 
\end{array}
\right).
$$
Then, by \eqref{cmc} we see that:
$$
(a_{ij})_{i,j=1,2}=
\left(
\begin{array}{cc}
b_{22}      &-b_{12}  \\
-b_{21}     &b_{11} 
\end{array}
\right)
$$
In particular, in this case the matrix 
$$
\frac1f(a_{ij})_{i,j=1,2}
$$ 
has the same diagonal entries of the matrix 
$$
\frac1f W_0,
$$ 
up to a permutation. The convexity of the $1$-homogeneous extension of the latter matrix then implies that assumption 
\eqref{cond-1 2} is verified, by the same argument used in the proof of Theorem \ref{main theorem n=2 rephrased}.
\end{proof}

\end{document}
